\textbf{1.} $a_1 = 2$, $a_2 = 3$, $a_3 = 7$, $a_4 = 25$, $a_5 = 121$. Les
trois premiers sont premiers ; $a_4 = 5^2$ et $a_5 = 11^2$ ne le sont pas.

\textbf{2.} Non : trois cas favorables ne prouvent rien, et le calcul
lui-même fournit des contre-exemples.

\textbf{3.} $a_4 = 25 = 5 \times 5$ n'est pas premier. Un seul contre-exemple
suffit à réfuter l'affirmation \og tous les termes sont premiers \fg.

\emph{Remarque.} D'autres termes sont premiers, par exemple
$a_{11} = 39\,916\,801$, mais on ne sait pas aujourd'hui s'il en existe une
infinité.
\hfill $\blacksquare$
